%%%%%%%%%%%%%%%%%%%%%%%%%%%%%%%%%%%%%%%%%%%%%%%%%%%%%%%%%%%%
% 14.11 SERIES EXPANSIONS
% The Composition of Advanced Mathematics
%%%%%%%%%%%%%%%%%%%%%%%%%%%%%%%%%%%%%%%%%%%%%%%%%%%%%%%%%%%%

\section{Series Expansions}
\label{sec:series-expansions}

\prerequisite{
Sequences and series, limits, derivatives, integrals,
exponential and logarithmic functions, trigonometric functions,
complex numbers, and basic convergence concepts.
}

\learningobjectives{
By the end of this reference section, the reader should be able to:
\begin{itemize}
    \item recognize finite and infinite series notation;
    \item distinguish sequences from series;
    \item recognize geometric series;
    \item use finite and infinite geometric-series formulas;
    \item recognize power series;
    \item identify centers and radii of convergence;
    \item recognize Taylor and Maclaurin series;
    \item construct Taylor polynomials;
    \item use standard series for exponential, logarithmic,
          trigonometric, inverse-trigonometric, and hyperbolic functions;
    \item recognize generalized binomial expansions;
    \item manipulate power series by addition, multiplication,
          differentiation, and integration;
    \item use series for local approximation;
    \item understand remainder and truncation error;
    \item recognize asymptotic expansions;
    \item distinguish convergent series from asymptotic series;
    \item recognize Fourier-series expansions conceptually;
    \item use complex exponential forms of Fourier series;
    \item connect generating functions with discrete mathematics;
    \item recognize common special-function expansions;
    \item avoid common errors involving convergence and termwise operations.
\end{itemize}
}

%%%%%%%%%%%%%%%%%%%%%%%%%%%%%%%%%%%%%%%%%%%%%%%%%%%%%%%%%%%%
\subsection{Sequences and Series}
\label{subsec:sequences-series-reference}
%%%%%%%%%%%%%%%%%%%%%%%%%%%%%%%%%%%%%%%%%%%%%%%%%%%%%%%%%%%%

A sequence is an ordered list

\[
\boxed{
a_1,a_2,a_3,\ldots
}
\]

or

\[
\boxed{
(a_n)_{n=1}^{\infty}.
}
\]

A series is a sum of sequence terms:

\[
\boxed{
\sum_{n=1}^{\infty}a_n.
}
\]

The partial sums are

\[
\boxed{
S_N
=
\sum_{n=1}^{N}a_n.
}
\]

The infinite series converges to \(S\) when

\[
\boxed{
\lim_{N\to\infty}S_N=S.
}
\]

%%%%%%%%%%%%%%%%%%%%%%%%%%%%%%%%%%%%%%%%%%%%%%%%%%%%%%%%%%%%
\subsection{Finite Geometric Series}
\label{subsec:finite-geometric-series-reference}
%%%%%%%%%%%%%%%%%%%%%%%%%%%%%%%%%%%%%%%%%%%%%%%%%%%%%%%%%%%%

For ratio \(r\neq1\),

\[
\boxed{
1+r+r^2+\cdots+r^{n}
=
\frac{1-r^{n+1}}{1-r}.
}
\]

More generally,

\[
\boxed{
\sum_{k=0}^{n}ar^k
=
a
\frac{1-r^{n+1}}{1-r}.
}
\]

%%%%%%%%%%%%%%%%%%%%%%%%%%%%%%%%%%%%%%%%%%%%%%%%%%%%%%%%%%%%
\subsection{Infinite Geometric Series}
\label{subsec:infinite-geometric-series-expansions}
%%%%%%%%%%%%%%%%%%%%%%%%%%%%%%%%%%%%%%%%%%%%%%%%%%%%%%%%%%%%

If

\[
|r|<1,
\]

then

\[
\boxed{
\sum_{n=0}^{\infty}r^n
=
\frac{1}{1-r}.
}
\]

More generally,

\[
\boxed{
\sum_{n=0}^{\infty}ar^n
=
\frac{a}{1-r}.
}
\]

If

\[
|r|\geq1,
\]

the geometric series does not converge.

%%%%%%%%%%%%%%%%%%%%%%%%%%%%%%%%%%%%%%%%%%%%%%%%%%%%%%%%%%%%
\subsection{Fundamental Geometric Expansion}
\label{subsec:fundamental-geometric-expansion}
%%%%%%%%%%%%%%%%%%%%%%%%%%%%%%%%%%%%%%%%%%%%%%%%%%%%%%%%%%%%

For

\[
|x|<1,
\]

\[
\boxed{
\frac{1}{1-x}
=
1+x+x^2+x^3+\cdots
}
\]

or

\[
\boxed{
\frac{1}{1-x}
=
\sum_{n=0}^{\infty}x^n.
}
\]

Replacing \(x\) by \(-x\) gives

\[
\boxed{
\frac{1}{1+x}
=
1-x+x^2-x^3+\cdots
}
\]

for \(|x|<1\).

%%%%%%%%%%%%%%%%%%%%%%%%%%%%%%%%%%%%%%%%%%%%%%%%%%%%%%%%%%%%
\subsection{Power Series}
\label{subsec:power-series-reference}
%%%%%%%%%%%%%%%%%%%%%%%%%%%%%%%%%%%%%%%%%%%%%%%%%%%%%%%%%%%%

A power series centered at \(a\) has form

\[
\boxed{
\sum_{n=0}^{\infty}
c_n(x-a)^n.
}
\]

The numbers \(c_n\) are coefficients.

The value \(a\) is the center.

%%%%%%%%%%%%%%%%%%%%%%%%%%%%%%%%%%%%%%%%%%%%%%%%%%%%%%%%%%%%
\subsection{Radius and Interval of Convergence}
\label{subsec:radius-interval-convergence-reference}
%%%%%%%%%%%%%%%%%%%%%%%%%%%%%%%%%%%%%%%%%%%%%%%%%%%%%%%%%%%%

A power series has a radius of convergence

\[
\boxed{
R\in[0,\infty].
}
\]

It converges absolutely when

\[
\boxed{
|x-a|<R,
}
\]

and diverges when

\[
\boxed{
|x-a|>R.
}
\]

Behavior at

\[
x=a\pm R
\]

must generally be tested separately.

%%%%%%%%%%%%%%%%%%%%%%%%%%%%%%%%%%%%%%%%%%%%%%%%%%%%%%%%%%%%
\subsection{Ratio Formula for Radius of Convergence}
\label{subsec:ratio-radius-convergence}
%%%%%%%%%%%%%%%%%%%%%%%%%%%%%%%%%%%%%%%%%%%%%%%%%%%%%%%%%%%%

If

\[
\lim_{n\to\infty}
\left|
\frac{c_n}{c_{n+1}}
\right|
=
R
\]

exists under suitable conditions, then this limit gives the radius of
convergence.

Equivalently, the ratio test may be applied directly to the power series.

%%%%%%%%%%%%%%%%%%%%%%%%%%%%%%%%%%%%%%%%%%%%%%%%%%%%%%%%%%%%
\subsection{Cauchy--Hadamard Formula}
\label{subsec:cauchy-hadamard-reference}
%%%%%%%%%%%%%%%%%%%%%%%%%%%%%%%%%%%%%%%%%%%%%%%%%%%%%%%%%%%%

For

\[
\sum_{n=0}^{\infty}c_n(x-a)^n,
\]

the radius of convergence satisfies

\[
\boxed{
\frac{1}{R}
=
\limsup_{n\to\infty}
|c_n|^{1/n}.
}
\]

This includes the conventions

\[
\frac{1}{0}=\infty,
\qquad
\frac{1}{\infty}=0.
\]

%%%%%%%%%%%%%%%%%%%%%%%%%%%%%%%%%%%%%%%%%%%%%%%%%%%%%%%%%%%%
\subsection{Taylor Series}
\label{subsec:taylor-series-reference}
%%%%%%%%%%%%%%%%%%%%%%%%%%%%%%%%%%%%%%%%%%%%%%%%%%%%%%%%%%%%

If \(f\) is sufficiently differentiable near \(a\), its Taylor series about
\(a\) is formally

\[
\boxed{
f(x)
=
\sum_{n=0}^{\infty}
\frac{
f^{(n)}(a)
}{
n!
}
(x-a)^n.
}
\]

Expanded,

\[
\boxed{
f(x)
=
f(a)
+
f'(a)(x-a)
+
\frac{f''(a)}{2!}(x-a)^2
+
\cdots.
}
\]

Whether the series actually equals \(f(x)\) depends on convergence and
remainder behavior.

%%%%%%%%%%%%%%%%%%%%%%%%%%%%%%%%%%%%%%%%%%%%%%%%%%%%%%%%%%%%
\subsection{Maclaurin Series}
\label{subsec:maclaurin-series-reference}
%%%%%%%%%%%%%%%%%%%%%%%%%%%%%%%%%%%%%%%%%%%%%%%%%%%%%%%%%%%%

A Taylor series centered at

\[
a=0
\]

is called a Maclaurin series:

\[
\boxed{
f(x)
=
\sum_{n=0}^{\infty}
\frac{
f^{(n)}(0)
}{
n!
}
x^n.
}
\]

%%%%%%%%%%%%%%%%%%%%%%%%%%%%%%%%%%%%%%%%%%%%%%%%%%%%%%%%%%%%
\subsection{Taylor Polynomial}
\label{subsec:taylor-polynomial-reference}
%%%%%%%%%%%%%%%%%%%%%%%%%%%%%%%%%%%%%%%%%%%%%%%%%%%%%%%%%%%%

The degree-\(N\) Taylor polynomial about \(a\) is

\[
\boxed{
T_N(x)
=
\sum_{n=0}^{N}
\frac{
f^{(n)}(a)
}{
n!
}
(x-a)^n.
}
\]

This provides a local polynomial approximation to \(f\).

%%%%%%%%%%%%%%%%%%%%%%%%%%%%%%%%%%%%%%%%%%%%%%%%%%%%%%%%%%%%
\subsection{Taylor Remainder}
\label{subsec:taylor-remainder-reference}
%%%%%%%%%%%%%%%%%%%%%%%%%%%%%%%%%%%%%%%%%%%%%%%%%%%%%%%%%%%%

The remainder is

\[
\boxed{
R_N(x)
=
f(x)-T_N(x).
}
\]

Under appropriate hypotheses, the Lagrange remainder has form

\[
\boxed{
R_N(x)
=
\frac{
f^{(N+1)}(\xi)
}{
(N+1)!
}
(x-a)^{N+1}
}
\]

for some \(\xi\) between \(a\) and \(x\).

%%%%%%%%%%%%%%%%%%%%%%%%%%%%%%%%%%%%%%%%%%%%%%%%%%%%%%%%%%%%
\subsection{Taylor Remainder Bound}
\label{subsec:taylor-remainder-bound}
%%%%%%%%%%%%%%%%%%%%%%%%%%%%%%%%%%%%%%%%%%%%%%%%%%%%%%%%%%%%

If

\[
|f^{(N+1)}(t)|
\leq M
\]

between \(a\) and \(x\), then

\[
\boxed{
|R_N(x)|
\leq
\frac{
M|x-a|^{N+1}
}{
(N+1)!
}.
}
\]

%%%%%%%%%%%%%%%%%%%%%%%%%%%%%%%%%%%%%%%%%%%%%%%%%%%%%%%%%%%%
\subsection{Exponential Series}
\label{subsec:exponential-series-reference}
%%%%%%%%%%%%%%%%%%%%%%%%%%%%%%%%%%%%%%%%%%%%%%%%%%%%%%%%%%%%

For every real or complex \(x\),

\[
\boxed{
e^x
=
\sum_{n=0}^{\infty}
\frac{x^n}{n!}.
}
\]

Thus

\[
\boxed{
e^x
=
1+x+\frac{x^2}{2!}
+\frac{x^3}{3!}
+\cdots.
}
\]

The radius of convergence is

\[
\boxed{
R=\infty.
}
\]

%%%%%%%%%%%%%%%%%%%%%%%%%%%%%%%%%%%%%%%%%%%%%%%%%%%%%%%%%%%%
\subsection{Worked Example: Approximate \(e\)}
\label{subsec:worked-approximate-e-series}
%%%%%%%%%%%%%%%%%%%%%%%%%%%%%%%%%%%%%%%%%%%%%%%%%%%%%%%%%%%%

\begin{workedexamplebox}{Maclaurin Approximation}

Use the first five terms of

\[
e^x
=
1+x+\frac{x^2}{2!}
+\frac{x^3}{3!}
+\frac{x^4}{4!}
+\cdots
\]

at \(x=1\).

Then

\[
e
\approx
1+1+\frac12+\frac16+\frac1{24}.
\]

Thus

\begin{answerbox}
\[
\boxed{
e
\approx
2.70833.
}
\]
\end{answerbox}

\end{workedexamplebox}

%%%%%%%%%%%%%%%%%%%%%%%%%%%%%%%%%%%%%%%%%%%%%%%%%%%%%%%%%%%%
\subsection{Sine Series}
\label{subsec:sine-series-reference}
%%%%%%%%%%%%%%%%%%%%%%%%%%%%%%%%%%%%%%%%%%%%%%%%%%%%%%%%%%%%

For every real or complex \(x\),

\[
\boxed{
\sin x
=
\sum_{n=0}^{\infty}
(-1)^n
\frac{
x^{2n+1}
}{
(2n+1)!
}.
}
\]

Thus

\[
\boxed{
\sin x
=
x
-
\frac{x^3}{3!}
+
\frac{x^5}{5!}
-
\frac{x^7}{7!}
+\cdots.
}
\]

%%%%%%%%%%%%%%%%%%%%%%%%%%%%%%%%%%%%%%%%%%%%%%%%%%%%%%%%%%%%
\subsection{Cosine Series}
\label{subsec:cosine-series-reference}
%%%%%%%%%%%%%%%%%%%%%%%%%%%%%%%%%%%%%%%%%%%%%%%%%%%%%%%%%%%%

\[
\boxed{
\cos x
=
\sum_{n=0}^{\infty}
(-1)^n
\frac{
x^{2n}
}{
(2n)!
}.
}
\]

Thus

\[
\boxed{
\cos x
=
1
-
\frac{x^2}{2!}
+
\frac{x^4}{4!}
-
\frac{x^6}{6!}
+\cdots.
}
\]

%%%%%%%%%%%%%%%%%%%%%%%%%%%%%%%%%%%%%%%%%%%%%%%%%%%%%%%%%%%%
\subsection{Euler Formula From Series}
\label{subsec:euler-formula-series-reference}
%%%%%%%%%%%%%%%%%%%%%%%%%%%%%%%%%%%%%%%%%%%%%%%%%%%%%%%%%%%%

Using

\[
e^{ix}
=
\sum_{n=0}^{\infty}
\frac{
(ix)^n
}{
n!
},
\]

separating even and odd powers yields

\[
\boxed{
e^{ix}
=
\cos x+i\sin x.
}
\]

%%%%%%%%%%%%%%%%%%%%%%%%%%%%%%%%%%%%%%%%%%%%%%%%%%%%%%%%%%%%
\subsection{Hyperbolic Cosine Series}
\label{subsec:cosh-series-reference}
%%%%%%%%%%%%%%%%%%%%%%%%%%%%%%%%%%%%%%%%%%%%%%%%%%%%%%%%%%%%

\[
\boxed{
\cosh x
=
\sum_{n=0}^{\infty}
\frac{
x^{2n}
}{
(2n)!
}.
}
\]

Thus

\[
\boxed{
\cosh x
=
1+\frac{x^2}{2!}
+\frac{x^4}{4!}
+\cdots.
}
\]

%%%%%%%%%%%%%%%%%%%%%%%%%%%%%%%%%%%%%%%%%%%%%%%%%%%%%%%%%%%%
\subsection{Hyperbolic Sine Series}
\label{subsec:sinh-series-reference}
%%%%%%%%%%%%%%%%%%%%%%%%%%%%%%%%%%%%%%%%%%%%%%%%%%%%%%%%%%%%

\[
\boxed{
\sinh x
=
\sum_{n=0}^{\infty}
\frac{
x^{2n+1}
}{
(2n+1)!
}.
}
\]

Thus

\[
\boxed{
\sinh x
=
x+\frac{x^3}{3!}
+\frac{x^5}{5!}
+\cdots.
}
\]

%%%%%%%%%%%%%%%%%%%%%%%%%%%%%%%%%%%%%%%%%%%%%%%%%%%%%%%%%%%%
\subsection{Logarithm Series}
\label{subsec:logarithm-series-reference}
%%%%%%%%%%%%%%%%%%%%%%%%%%%%%%%%%%%%%%%%%%%%%%%%%%%%%%%%%%%%

For

\[
-1<x\leq1,
\]

with special endpoint interpretation,

\[
\boxed{
\ln(1+x)
=
\sum_{n=1}^{\infty}
(-1)^{n+1}
\frac{x^n}{n}.
}
\]

Thus

\[
\boxed{
\ln(1+x)
=
x
-\frac{x^2}{2}
+\frac{x^3}{3}
-\frac{x^4}{4}
+\cdots.
}
\]

The power-series radius is

\[
\boxed{
R=1.
}
\]

%%%%%%%%%%%%%%%%%%%%%%%%%%%%%%%%%%%%%%%%%%%%%%%%%%%%%%%%%%%%
\subsection{Series for \(-\ln(1-x)\)}
\label{subsec:negative-log-one-minus-x}
%%%%%%%%%%%%%%%%%%%%%%%%%%%%%%%%%%%%%%%%%%%%%%%%%%%%%%%%%%%%

For

\[
|x|<1,
\]

\[
\boxed{
-\ln(1-x)
=
\sum_{n=1}^{\infty}
\frac{x^n}{n}.
}
\]

Thus

\[
\boxed{
-\ln(1-x)
=
x+\frac{x^2}{2}
+\frac{x^3}{3}
+\cdots.
}
\]

%%%%%%%%%%%%%%%%%%%%%%%%%%%%%%%%%%%%%%%%%%%%%%%%%%%%%%%%%%%%
\subsection{Arctangent Series}
\label{subsec:arctangent-series-reference}
%%%%%%%%%%%%%%%%%%%%%%%%%%%%%%%%%%%%%%%%%%%%%%%%%%%%%%%%%%%%

For

\[
|x|\leq1
\]

with endpoint convergence interpreted appropriately,

\[
\boxed{
\arctan x
=
\sum_{n=0}^{\infty}
(-1)^n
\frac{
x^{2n+1}
}{
2n+1
}.
}
\]

Thus

\[
\boxed{
\arctan x
=
x-\frac{x^3}{3}
+\frac{x^5}{5}
-\frac{x^7}{7}
+\cdots.
}
\]

%%%%%%%%%%%%%%%%%%%%%%%%%%%%%%%%%%%%%%%%%%%%%%%%%%%%%%%%%%%%
\subsection{Gregory--Leibniz Series}
\label{subsec:gregory-leibniz-series-reference}
%%%%%%%%%%%%%%%%%%%%%%%%%%%%%%%%%%%%%%%%%%%%%%%%%%%%%%%%%%%%

Set

\[
x=1
\]

in the arctangent series:

\[
\frac{\pi}{4}
=
1-\frac13+\frac15-\frac17+\cdots.
\]

Therefore,

\[
\boxed{
\pi
=
4
\sum_{n=0}^{\infty}
\frac{
(-1)^n
}{
2n+1
}.
}
\]

%%%%%%%%%%%%%%%%%%%%%%%%%%%%%%%%%%%%%%%%%%%%%%%%%%%%%%%%%%%%
\subsection{Arcsine Series}
\label{subsec:arcsine-series-reference}
%%%%%%%%%%%%%%%%%%%%%%%%%%%%%%%%%%%%%%%%%%%%%%%%%%%%%%%%%%%%

For \(|x|<1\),

\[
\boxed{
\arcsin x
=
\sum_{n=0}^{\infty}
\frac{
(2n)!
}{
4^n(n!)^2(2n+1)
}
x^{2n+1}.
}
\]

The first terms are

\[
\boxed{
\arcsin x
=
x
+
\frac{x^3}{6}
+
\frac{3x^5}{40}
+
\frac{5x^7}{112}
+\cdots.
}
\]

%%%%%%%%%%%%%%%%%%%%%%%%%%%%%%%%%%%%%%%%%%%%%%%%%%%%%%%%%%%%
\subsection{Binomial Series}
\label{subsec:binomial-series-reference}
%%%%%%%%%%%%%%%%%%%%%%%%%%%%%%%%%%%%%%%%%%%%%%%%%%%%%%%%%%%%

For general exponent \(\alpha\),

\[
\boxed{
(1+x)^\alpha
=
\sum_{n=0}^{\infty}
\binom{\alpha}{n}
x^n
}
\]

for \(|x|<1\), where

\[
\boxed{
\binom{\alpha}{n}
=
\frac{
\alpha(\alpha-1)\cdots(\alpha-n+1)
}{
n!
}.
}
\]

%%%%%%%%%%%%%%%%%%%%%%%%%%%%%%%%%%%%%%%%%%%%%%%%%%%%%%%%%%%%
\subsection{First Terms of the Binomial Series}
\label{subsec:binomial-series-first-terms}
%%%%%%%%%%%%%%%%%%%%%%%%%%%%%%%%%%%%%%%%%%%%%%%%%%%%%%%%%%%%

\[
\boxed{
(1+x)^\alpha
=
1
+
\alpha x
+
\frac{
\alpha(\alpha-1)
}{
2!
}
x^2
+
\frac{
\alpha(\alpha-1)(\alpha-2)
}{
3!
}
x^3
+\cdots.
}
\]

%%%%%%%%%%%%%%%%%%%%%%%%%%%%%%%%%%%%%%%%%%%%%%%%%%%%%%%%%%%%
\subsection{Square-Root Expansion}
\label{subsec:sqrt-series-reference}
%%%%%%%%%%%%%%%%%%%%%%%%%%%%%%%%%%%%%%%%%%%%%%%%%%%%%%%%%%%%

Using

\[
\alpha=\frac12,
\]

\[
\boxed{
\sqrt{1+x}
=
1
+
\frac{x}{2}
-
\frac{x^2}{8}
+
\frac{x^3}{16}
-
\frac{5x^4}{128}
+\cdots
}
\]

for \(|x|<1\).

%%%%%%%%%%%%%%%%%%%%%%%%%%%%%%%%%%%%%%%%%%%%%%%%%%%%%%%%%%%%
\subsection{Reciprocal Square-Root Expansion}
\label{subsec:inverse-sqrt-series-reference}
%%%%%%%%%%%%%%%%%%%%%%%%%%%%%%%%%%%%%%%%%%%%%%%%%%%%%%%%%%%%

Using

\[
\alpha=-\frac12,
\]

\[
\boxed{
\frac{1}{\sqrt{1+x}}
=
1
-
\frac{x}{2}
+
\frac{3x^2}{8}
-
\frac{5x^3}{16}
+
\frac{35x^4}{128}
-\cdots.
}
\]

%%%%%%%%%%%%%%%%%%%%%%%%%%%%%%%%%%%%%%%%%%%%%%%%%%%%%%%%%%%%
\subsection{Reciprocal Expansion}
\label{subsec:reciprocal-series-reference}
%%%%%%%%%%%%%%%%%%%%%%%%%%%%%%%%%%%%%%%%%%%%%%%%%%%%%%%%%%%%

For \(|x|<1\),

\[
\boxed{
\frac{1}{1+x}
=
1-x+x^2-x^3+\cdots.
}
\]

Also,

\[
\boxed{
\frac{1}{(1-x)^2}
=
\sum_{n=0}^{\infty}
(n+1)x^n.
}
\]

%%%%%%%%%%%%%%%%%%%%%%%%%%%%%%%%%%%%%%%%%%%%%%%%%%%%%%%%%%%%
\subsection{Differentiating Power Series}
\label{subsec:differentiate-power-series-reference}
%%%%%%%%%%%%%%%%%%%%%%%%%%%%%%%%%%%%%%%%%%%%%%%%%%%%%%%%%%%%

Inside its radius of convergence,

\[
\sum_{n=0}^{\infty}
c_n(x-a)^n
\]

may be differentiated term by term:

\[
\boxed{
\frac{d}{dx}
\sum_{n=0}^{\infty}
c_n(x-a)^n
=
\sum_{n=1}^{\infty}
nc_n(x-a)^{n-1}.
}
\]

The differentiated series has the same radius of convergence.

%%%%%%%%%%%%%%%%%%%%%%%%%%%%%%%%%%%%%%%%%%%%%%%%%%%%%%%%%%%%
\subsection{Integrating Power Series}
\label{subsec:integrate-power-series-reference}
%%%%%%%%%%%%%%%%%%%%%%%%%%%%%%%%%%%%%%%%%%%%%%%%%%%%%%%%%%%%

Inside the radius of convergence,

\[
\boxed{
\int
\left[
\sum_{n=0}^{\infty}
c_n(x-a)^n
\right]
dx
=
C
+
\sum_{n=0}^{\infty}
\frac{
c_n
}{
n+1
}
(x-a)^{n+1}.
}
\]

The integrated series has the same radius of convergence.

%%%%%%%%%%%%%%%%%%%%%%%%%%%%%%%%%%%%%%%%%%%%%%%%%%%%%%%%%%%%
\subsection{Worked Example: Derive a Series by Integration}
\label{subsec:worked-series-integration}
%%%%%%%%%%%%%%%%%%%%%%%%%%%%%%%%%%%%%%%%%%%%%%%%%%%%%%%%%%%%

\begin{workedexamplebox}{Derive the Logarithm Series}

Start with

\[
\frac{1}{1+x}
=
1-x+x^2-x^3+\cdots,
\qquad
|x|<1.
\]

Integrate from \(0\) to \(x\):

\[
\int_0^x
\frac{1}{1+t}
\,dt
=
\int_0^x
(1-t+t^2-t^3+\cdots)
\,dt.
\]

Thus

\[
\ln(1+x)
=
x-\frac{x^2}{2}
+\frac{x^3}{3}
-\frac{x^4}{4}
+\cdots.
\]

\begin{answerbox}
\[
\boxed{
\ln(1+x)
=
\sum_{n=1}^{\infty}
(-1)^{n+1}
\frac{x^n}{n},
\qquad
|x|<1.
}
\]
\end{answerbox}

\end{workedexamplebox}

%%%%%%%%%%%%%%%%%%%%%%%%%%%%%%%%%%%%%%%%%%%%%%%%%%%%%%%%%%%%
\subsection{Multiplication of Power Series}
\label{subsec:multiply-power-series-reference}
%%%%%%%%%%%%%%%%%%%%%%%%%%%%%%%%%%%%%%%%%%%%%%%%%%%%%%%%%%%%

If

\[
A(x)
=
\sum_{n=0}^{\infty}a_nx^n
\]

and

\[
B(x)
=
\sum_{n=0}^{\infty}b_nx^n,
\]

then their product is given by the Cauchy product

\[
\boxed{
A(x)B(x)
=
\sum_{n=0}^{\infty}
\left(
\sum_{k=0}^{n}
a_kb_{n-k}
\right)
x^n
}
\]

under appropriate convergence conditions.

%%%%%%%%%%%%%%%%%%%%%%%%%%%%%%%%%%%%%%%%%%%%%%%%%%%%%%%%%%%%
\subsection{Composition of Series}
\label{subsec:composition-series-reference}
%%%%%%%%%%%%%%%%%%%%%%%%%%%%%%%%%%%%%%%%%%%%%%%%%%%%%%%%%%%%

Power series may often be composed by replacing \(x\) with another expression.

For example,

\[
e^x
=
1+x+\frac{x^2}{2!}+\cdots
\]

implies

\[
\boxed{
e^{-x^2}
=
1
-
x^2
+
\frac{x^4}{2!}
-
\frac{x^6}{3!}
+\cdots.
}
\]

%%%%%%%%%%%%%%%%%%%%%%%%%%%%%%%%%%%%%%%%%%%%%%%%%%%%%%%%%%%%
\subsection{Worked Example: Local Approximation}
\label{subsec:worked-local-series-approximation}
%%%%%%%%%%%%%%%%%%%%%%%%%%%%%%%%%%%%%%%%%%%%%%%%%%%%%%%%%%%%

\begin{workedexamplebox}{Approximate \(\sin(0.1)\)}

Use

\[
\sin x
=
x-\frac{x^3}{6}
+\frac{x^5}{120}
-\cdots.
\]

At

\[
x=0.1,
\]

using the first two nonzero terms,

\[
\sin(0.1)
\approx
0.1
-
\frac{0.1^3}{6}.
\]

Therefore,

\begin{answerbox}
\[
\boxed{
\sin(0.1)
\approx
0.0998333.
}
\]
\end{answerbox}

\end{workedexamplebox}

%%%%%%%%%%%%%%%%%%%%%%%%%%%%%%%%%%%%%%%%%%%%%%%%%%%%%%%%%%%%
\subsection{Alternating-Series Error Estimate}
\label{subsec:alternating-series-error-reference}
%%%%%%%%%%%%%%%%%%%%%%%%%%%%%%%%%%%%%%%%%%%%%%%%%%%%%%%%%%%%

For an alternating series satisfying the standard alternating-series
conditions, the truncation error after \(N\) terms satisfies

\[
\boxed{
|R_N|
\leq
\text{magnitude of the first omitted term}.
}
\]

This is useful for series such as sine, cosine, and arctangent in suitable
regions.

%%%%%%%%%%%%%%%%%%%%%%%%%%%%%%%%%%%%%%%%%%%%%%%%%%%%%%%%%%%%
\subsection{Series-Based Limits}
\label{subsec:series-based-limits-reference}
%%%%%%%%%%%%%%%%%%%%%%%%%%%%%%%%%%%%%%%%%%%%%%%%%%%%%%%%%%%%

Series can expose leading-order behavior.

For example,

\[
\sin x
=
x-\frac{x^3}{6}+O(x^5),
\]

so

\[
\sin x-x
=
-\frac{x^3}{6}+O(x^5).
\]

Therefore,

\[
\boxed{
\lim_{x\to0}
\frac{
\sin x-x
}{
x^3
}
=
-\frac16.
}
\]

%%%%%%%%%%%%%%%%%%%%%%%%%%%%%%%%%%%%%%%%%%%%%%%%%%%%%%%%%%%%
\subsection{Common Local Expansions}
\label{subsec:common-local-expansions}
%%%%%%%%%%%%%%%%%%%%%%%%%%%%%%%%%%%%%%%%%%%%%%%%%%%%%%%%%%%%

As \(x\to0\),

\[
\boxed{
e^x
=
1+x+\frac{x^2}{2}
+\frac{x^3}{6}
+O(x^4),
}
\]

\[
\boxed{
\sin x
=
x-\frac{x^3}{6}
+O(x^5),
}
\]

\[
\boxed{
\cos x
=
1-\frac{x^2}{2}
+\frac{x^4}{24}
+O(x^6),
}
\]

\[
\boxed{
\tan x
=
x+\frac{x^3}{3}
+O(x^5),
}
\]

\[
\boxed{
\ln(1+x)
=
x-\frac{x^2}{2}
+\frac{x^3}{3}
+O(x^4),
}
\]

and

\[
\boxed{
(1+x)^\alpha
=
1+\alpha x
+
\frac{\alpha(\alpha-1)}{2}x^2
+
O(x^3).
}
\]

%%%%%%%%%%%%%%%%%%%%%%%%%%%%%%%%%%%%%%%%%%%%%%%%%%%%%%%%%%%%
\subsection{Asymptotic Expansions}
\label{subsec:asymptotic-expansions-reference}
%%%%%%%%%%%%%%%%%%%%%%%%%%%%%%%%%%%%%%%%%%%%%%%%%%%%%%%%%%%%

An asymptotic expansion

\[
\boxed{
f(x)
\sim
a_0\phi_0(x)
+
a_1\phi_1(x)
+
a_2\phi_2(x)
+\cdots
}
\]

means successive truncations approximate \(f(x)\) to progressively smaller
orders in the specified limiting regime.

An asymptotic series need not converge as an infinite series.

%%%%%%%%%%%%%%%%%%%%%%%%%%%%%%%%%%%%%%%%%%%%%%%%%%%%%%%%%%%%
\subsection{Convergent Series Versus Asymptotic Series}
\label{subsec:convergent-vs-asymptotic-series}
%%%%%%%%%%%%%%%%%%%%%%%%%%%%%%%%%%%%%%%%%%%%%%%%%%%%%%%%%%%%

A convergent power series satisfies

\[
\boxed{
f(x)
=
\sum_{n=0}^{\infty}a_n(x-a)^n
}
\]

within a region of convergence.

An asymptotic expansion may instead satisfy

\[
\boxed{
f(x)
-
\sum_{n=0}^{N}
a_n\phi_n(x)
=
o(\phi_N(x))
}
\]

for every fixed \(N\), without the infinite series converging.

%%%%%%%%%%%%%%%%%%%%%%%%%%%%%%%%%%%%%%%%%%%%%%%%%%%%%%%%%%%%
\subsection{Stirling Series}
\label{subsec:stirling-series-reference}
%%%%%%%%%%%%%%%%%%%%%%%%%%%%%%%%%%%%%%%%%%%%%%%%%%%%%%%%%%%%

As \(n\to\infty\),

\[
\boxed{
n!
\sim
\sqrt{2\pi n}
\left(
\frac{n}{e}
\right)^n.
}
\]

A refined expansion is

\[
\boxed{
n!
\sim
\sqrt{2\pi n}
\left(
\frac{n}{e}
\right)^n
\left[
1
+
\frac{1}{12n}
+
\frac{1}{288n^2}
-
\frac{139}{51840n^3}
+\cdots
\right].
}
\]

%%%%%%%%%%%%%%%%%%%%%%%%%%%%%%%%%%%%%%%%%%%%%%%%%%%%%%%%%%%%
\subsection{Harmonic-Number Expansion}
\label{subsec:harmonic-number-expansion-reference}
%%%%%%%%%%%%%%%%%%%%%%%%%%%%%%%%%%%%%%%%%%%%%%%%%%%%%%%%%%%%

As \(n\to\infty\),

\[
\boxed{
H_n
=
\ln n
+
\gamma
+
\frac{1}{2n}
-
\frac{1}{12n^2}
+
\frac{1}{120n^4}
+\cdots.
}
\]

%%%%%%%%%%%%%%%%%%%%%%%%%%%%%%%%%%%%%%%%%%%%%%%%%%%%%%%%%%%%
\subsection{Fourier Series}
\label{subsec:fourier-series-reference}
%%%%%%%%%%%%%%%%%%%%%%%%%%%%%%%%%%%%%%%%%%%%%%%%%%%%%%%%%%%%

A \(2\pi\)-periodic function may, under suitable conditions, be represented
by

\[
\boxed{
f(x)
\sim
\frac{a_0}{2}
+
\sum_{n=1}^{\infty}
\left(
a_n\cos nx
+
b_n\sin nx
\right).
}
\]

The Fourier coefficients are

\[
\boxed{
a_n
=
\frac{1}{\pi}
\int_{-\pi}^{\pi}
f(x)\cos(nx)\,dx,
}
\]

and

\[
\boxed{
b_n
=
\frac{1}{\pi}
\int_{-\pi}^{\pi}
f(x)\sin(nx)\,dx.
}
\]

Also,

\[
\boxed{
a_0
=
\frac{1}{\pi}
\int_{-\pi}^{\pi}
f(x)\,dx.
}
\]

%%%%%%%%%%%%%%%%%%%%%%%%%%%%%%%%%%%%%%%%%%%%%%%%%%%%%%%%%%%%
\subsection{Fourier Series on \([-L,L]\)}
\label{subsec:fourier-series-general-interval}
%%%%%%%%%%%%%%%%%%%%%%%%%%%%%%%%%%%%%%%%%%%%%%%%%%%%%%%%%%%%

For period \(2L\),

\[
\boxed{
f(x)
\sim
\frac{a_0}{2}
+
\sum_{n=1}^{\infty}
\left[
a_n
\cos\left(
\frac{n\pi x}{L}
\right)
+
b_n
\sin\left(
\frac{n\pi x}{L}
\right)
\right].
}
\]

Here

\[
\boxed{
a_n
=
\frac1L
\int_{-L}^{L}
f(x)
\cos\left(
\frac{n\pi x}{L}
\right)
dx,
}
\]

and

\[
\boxed{
b_n
=
\frac1L
\int_{-L}^{L}
f(x)
\sin\left(
\frac{n\pi x}{L}
\right)
dx.
}
\]

%%%%%%%%%%%%%%%%%%%%%%%%%%%%%%%%%%%%%%%%%%%%%%%%%%%%%%%%%%%%
\subsection{Even Fourier Series}
\label{subsec:even-fourier-series-reference}
%%%%%%%%%%%%%%%%%%%%%%%%%%%%%%%%%%%%%%%%%%%%%%%%%%%%%%%%%%%%

If \(f\) is even, then

\[
\boxed{
b_n=0.
}
\]

Therefore the Fourier series contains only cosine terms:

\[
\boxed{
f(x)
\sim
\frac{a_0}{2}
+
\sum_{n=1}^{\infty}
a_n\cos(nx).
}
\]

%%%%%%%%%%%%%%%%%%%%%%%%%%%%%%%%%%%%%%%%%%%%%%%%%%%%%%%%%%%%
\subsection{Odd Fourier Series}
\label{subsec:odd-fourier-series-reference}
%%%%%%%%%%%%%%%%%%%%%%%%%%%%%%%%%%%%%%%%%%%%%%%%%%%%%%%%%%%%

If \(f\) is odd, then

\[
\boxed{
a_n=0
}
\]

for all \(n\geq0\).

Thus

\[
\boxed{
f(x)
\sim
\sum_{n=1}^{\infty}
b_n\sin(nx).
}
\]

%%%%%%%%%%%%%%%%%%%%%%%%%%%%%%%%%%%%%%%%%%%%%%%%%%%%%%%%%%%%
\subsection{Complex Fourier Series}
\label{subsec:complex-fourier-series-reference}
%%%%%%%%%%%%%%%%%%%%%%%%%%%%%%%%%%%%%%%%%%%%%%%%%%%%%%%%%%%%

A \(2\pi\)-periodic function may be represented as

\[
\boxed{
f(x)
\sim
\sum_{n=-\infty}^{\infty}
c_n e^{inx},
}
\]

where

\[
\boxed{
c_n
=
\frac{1}{2\pi}
\int_{-\pi}^{\pi}
f(x)e^{-inx}\,dx.
}
\]

%%%%%%%%%%%%%%%%%%%%%%%%%%%%%%%%%%%%%%%%%%%%%%%%%%%%%%%%%%%%
\subsection{Parseval Identity}
\label{subsec:parseval-series-reference}
%%%%%%%%%%%%%%%%%%%%%%%%%%%%%%%%%%%%%%%%%%%%%%%%%%%%%%%%%%%%

Under suitable hypotheses,

\[
\boxed{
\frac1\pi
\int_{-\pi}^{\pi}
|f(x)|^2
\,dx
=
\frac{a_0^2}{2}
+
\sum_{n=1}^{\infty}
(a_n^2+b_n^2).
}
\]

In complex form,

\[
\boxed{
\frac{1}{2\pi}
\int_{-\pi}^{\pi}
|f(x)|^2
\,dx
=
\sum_{n=-\infty}^{\infty}
|c_n|^2.
}
\]

%%%%%%%%%%%%%%%%%%%%%%%%%%%%%%%%%%%%%%%%%%%%%%%%%%%%%%%%%%%%
\subsection{Generating Functions}
\label{subsec:generating-functions-reference}
%%%%%%%%%%%%%%%%%%%%%%%%%%%%%%%%%%%%%%%%%%%%%%%%%%%%%%%%%%%%

For a sequence \((a_n)\), its ordinary generating function is

\[
\boxed{
A(x)
=
\sum_{n=0}^{\infty}
a_nx^n.
}
\]

Generating functions encode discrete sequences as formal or analytic power
series.

%%%%%%%%%%%%%%%%%%%%%%%%%%%%%%%%%%%%%%%%%%%%%%%%%%%%%%%%%%%%
\subsection{Generating Function for a Constant Sequence}
\label{subsec:generating-function-constant-sequence}
%%%%%%%%%%%%%%%%%%%%%%%%%%%%%%%%%%%%%%%%%%%%%%%%%%%%%%%%%%%%

For

\[
a_n=1,
\]

\[
\boxed{
A(x)
=
\sum_{n=0}^{\infty}
x^n
=
\frac{1}{1-x}.
}
\]

%%%%%%%%%%%%%%%%%%%%%%%%%%%%%%%%%%%%%%%%%%%%%%%%%%%%%%%%%%%%
\subsection{Generating Function for Natural Numbers}
\label{subsec:generating-function-natural-numbers}
%%%%%%%%%%%%%%%%%%%%%%%%%%%%%%%%%%%%%%%%%%%%%%%%%%%%%%%%%%%%

Differentiate

\[
\frac{1}{1-x}
=
\sum_{n=0}^{\infty}x^n.
\]

Then

\[
\boxed{
\frac{1}{(1-x)^2}
=
\sum_{n=1}^{\infty}
nx^{n-1}.
}
\]

Multiplying by \(x\),

\[
\boxed{
\sum_{n=1}^{\infty}
nx^n
=
\frac{x}{(1-x)^2}.
}
\]

%%%%%%%%%%%%%%%%%%%%%%%%%%%%%%%%%%%%%%%%%%%%%%%%%%%%%%%%%%%%
\subsection{Exponential Generating Functions}
\label{subsec:exponential-generating-functions-reference}
%%%%%%%%%%%%%%%%%%%%%%%%%%%%%%%%%%%%%%%%%%%%%%%%%%%%%%%%%%%%

The exponential generating function of \((a_n)\) is

\[
\boxed{
A(x)
=
\sum_{n=0}^{\infty}
a_n
\frac{x^n}{n!}.
}
\]

These are especially useful in combinatorics.

%%%%%%%%%%%%%%%%%%%%%%%%%%%%%%%%%%%%%%%%%%%%%%%%%%%%%%%%%%%%
\subsection{Special-Function Series}
\label{subsec:special-function-series-reference}
%%%%%%%%%%%%%%%%%%%%%%%%%%%%%%%%%%%%%%%%%%%%%%%%%%%%%%%%%%%%

Many special functions are defined or represented by infinite series.

For example, the Bessel function \(J_0\) has expansion

\[
\boxed{
J_0(x)
=
\sum_{n=0}^{\infty}
(-1)^n
\frac{
(x^2/4)^n
}{
(n!)^2
}.
}
\]

Equivalently,

\[
\boxed{
J_0(x)
=
1
-
\frac{x^2}{4}
+
\frac{x^4}{64}
-\cdots.
}
\]

%%%%%%%%%%%%%%%%%%%%%%%%%%%%%%%%%%%%%%%%%%%%%%%%%%%%%%%%%%%%
\subsection{Error Function Series}
\label{subsec:error-function-series-reference}
%%%%%%%%%%%%%%%%%%%%%%%%%%%%%%%%%%%%%%%%%%%%%%%%%%%%%%%%%%%%

The error function is

\[
\boxed{
\operatorname{erf}(x)
=
\frac{2}{\sqrt{\pi}}
\int_0^x
e^{-t^2}
\,dt.
}
\]

Using the exponential series,

\[
\boxed{
\operatorname{erf}(x)
=
\frac{2}{\sqrt{\pi}}
\sum_{n=0}^{\infty}
\frac{
(-1)^n x^{2n+1}
}{
n!(2n+1)
}.
}
\]

%%%%%%%%%%%%%%%%%%%%%%%%%%%%%%%%%%%%%%%%%%%%%%%%%%%%%%%%%%%%
\subsection{Quick Maclaurin Table}
\label{subsec:quick-maclaurin-table}
%%%%%%%%%%%%%%%%%%%%%%%%%%%%%%%%%%%%%%%%%%%%%%%%%%%%%%%%%%%%

\[
\boxed{
e^x
=
1+x+\frac{x^2}{2!}
+\frac{x^3}{3!}
+\cdots
}
\]

\[
\boxed{
\sin x
=
x-\frac{x^3}{3!}
+\frac{x^5}{5!}
-\cdots
}
\]

\[
\boxed{
\cos x
=
1-\frac{x^2}{2!}
+\frac{x^4}{4!}
-\cdots
}
\]

\[
\boxed{
\frac{1}{1-x}
=
1+x+x^2+x^3+\cdots
}
\]

\[
\boxed{
\ln(1+x)
=
x-\frac{x^2}{2}
+\frac{x^3}{3}
-\cdots
}
\]

\[
\boxed{
\arctan x
=
x-\frac{x^3}{3}
+\frac{x^5}{5}
-\cdots
}
\]

\[
\boxed{
\sinh x
=
x+\frac{x^3}{3!}
+\frac{x^5}{5!}
+\cdots
}
\]

\[
\boxed{
\cosh x
=
1+\frac{x^2}{2!}
+\frac{x^4}{4!}
+\cdots.
}
\]

%%%%%%%%%%%%%%%%%%%%%%%%%%%%%%%%%%%%%%%%%%%%%%%%%%%%%%%%%%%%
\subsection{Common Errors}
\label{subsec:series-expansions-common-errors}
%%%%%%%%%%%%%%%%%%%%%%%%%%%%%%%%%%%%%%%%%%%%%%%%%%%%%%%%%%%%

\subsubsection{Confusing a Sequence With a Series}

\[
a_n
\]

is a sequence term.

\[
\sum a_n
\]

is a series.

Their convergence properties are related but distinct.

%%%%%%%%%%%%%%%%%%%%%%%%%%%%%%%%%%%%%%%%%%%%%%%%%%%%%%%%%%%%
\subsubsection{Ignoring the Radius of Convergence}

A power-series identity may only hold within its convergence interval.

For example,

\[
\frac{1}{1-x}
=
\sum_{n=0}^{\infty}x^n
\]

requires

\[
\boxed{
|x|<1.
}
\]

%%%%%%%%%%%%%%%%%%%%%%%%%%%%%%%%%%%%%%%%%%%%%%%%%%%%%%%%%%%%
\subsubsection{Ignoring Endpoint Tests}

If a radius is

\[
R,
\]

then the points

\[
x=a-R,
\qquad
x=a+R
\]

usually require separate convergence analysis.

%%%%%%%%%%%%%%%%%%%%%%%%%%%%%%%%%%%%%%%%%%%%%%%%%%%%%%%%%%%%
\subsubsection{Assuming Taylor Series Always Equals the Function}

A function may be infinitely differentiable while its Taylor series fails to
equal the function away from the expansion point.

Analyticity is stronger than smoothness.

%%%%%%%%%%%%%%%%%%%%%%%%%%%%%%%%%%%%%%%%%%%%%%%%%%%%%%%%%%%%
\subsubsection{Dropping Factorials}

Taylor coefficients contain

\[
\boxed{
\frac{f^{(n)}(a)}{n!}.
}
\]

Forgetting \(n!\) produces incorrect expansions.

%%%%%%%%%%%%%%%%%%%%%%%%%%%%%%%%%%%%%%%%%%%%%%%%%%%%%%%%%%%%
\subsubsection{Using Degrees in Trigonometric Series}

The standard expansions

\[
\sin x
=
x-\frac{x^3}{6}+\cdots
\]

and

\[
\cos x
=
1-\frac{x^2}{2}+\cdots
\]

assume radians.

%%%%%%%%%%%%%%%%%%%%%%%%%%%%%%%%%%%%%%%%%%%%%%%%%%%%%%%%%%%%
\subsubsection{Confusing Equality With Asymptotic Equality}

The notation

\[
f(x)\sim g(x)
\]

does not mean

\[
f(x)=g(x).
\]

It means their ratio tends to \(1\) in the specified limit.

%%%%%%%%%%%%%%%%%%%%%%%%%%%%%%%%%%%%%%%%%%%%%%%%%%%%%%%%%%%%
\subsubsection{Assuming an Asymptotic Series Converges}

An asymptotic expansion can be useful even when its infinite series diverges.

Often an optimally truncated divergent series gives excellent approximation.

%%%%%%%%%%%%%%%%%%%%%%%%%%%%%%%%%%%%%%%%%%%%%%%%%%%%%%%%%%%%
\subsubsection{Differentiating Outside the Convergence Interval}

Termwise differentiation is justified inside the radius of convergence.

Boundary behavior requires additional analysis.

%%%%%%%%%%%%%%%%%%%%%%%%%%%%%%%%%%%%%%%%%%%%%%%%%%%%%%%%%%%%
\subsubsection{Ignoring Fourier Convergence Conditions}

A Fourier series may converge to the midpoint of a jump rather than to the
function's assigned point value.

Convergence depends on regularity conditions.

%%%%%%%%%%%%%%%%%%%%%%%%%%%%%%%%%%%%%%%%%%%%%%%%%%%%%%%%%%%%
\subsection{Series Expansion Workflow}
\label{subsec:series-expansion-workflow}
%%%%%%%%%%%%%%%%%%%%%%%%%%%%%%%%%%%%%%%%%%%%%%%%%%%%%%%%%%%%

A practical workflow is:

\begin{enumerate}
    \item identify whether the object is a numerical series or power series;
    \item locate a known standard expansion;
    \item shift or scale the expansion when possible;
    \item identify the center;
    \item determine the radius of convergence;
    \item test finite endpoints separately;
    \item differentiate or integrate termwise only where justified;
    \item use only as many terms as the desired approximation requires;
    \item estimate truncation error;
    \item distinguish exact convergent representations from asymptotic ones;
    \item verify leading terms by differentiation or limits when practical.
\end{enumerate}

%%%%%%%%%%%%%%%%%%%%%%%%%%%%%%%%%%%%%%%%%%%%%%%%%%%%%%%%%%%%
\subsection{Exercises}
\label{subsec:series-expansions-exercises}
%%%%%%%%%%%%%%%%%%%%%%%%%%%%%%%%%%%%%%%%%%%%%%%%%%%%%%%%%%%%

\begin{exercise}[1]
Write the first five terms of

\[
\sum_{n=0}^{\infty}x^n.
\]
\end{exercise}

\begin{exercise}[1]
Find the sum of

\[
1+\frac12+\frac14+\frac18+\cdots.
\]
\end{exercise}

\begin{exercise}[1]
State the convergence condition for a geometric series.
\end{exercise}

\begin{exercise}[1]
Write the Maclaurin series for \(e^x\).
\end{exercise}

\begin{exercise}[1]
Write the Maclaurin series for \(\sin x\).
\end{exercise}

\begin{exercise}[1]
Write the Maclaurin series for \(\cos x\).
\end{exercise}

\begin{exercise}[2]
Find the degree-4 Maclaurin polynomial for \(e^x\).
\end{exercise}

\begin{exercise}[2]
Find the degree-5 Maclaurin polynomial for \(\sin x\).
\end{exercise}

\begin{exercise}[2]
Find the degree-6 Maclaurin polynomial for \(\cos x\).
\end{exercise}

\begin{exercise}[2]
Approximate \(e^{0.2}\) using a Taylor polynomial.
\end{exercise}

\begin{exercise}[2]
Approximate \(\sin(0.2)\) using three nonzero terms.
\end{exercise}

\begin{exercise}[2]
Use the alternating-series remainder estimate for the previous approximation.
\end{exercise}

\begin{exercise}[2]
Derive the series for

\[
\frac{1}{1+x}
\]

from the geometric series.
\end{exercise}

\begin{exercise}[2]
Derive the series for

\[
\frac{1}{(1-x)^2}
\]

by differentiation.
\end{exercise}

\begin{exercise}[2]
Derive the series for

\[
-\ln(1-x)
\]

by integration.
\end{exercise}

\begin{exercise}[2]
Find the series for

\[
e^{-x}.
\]
\end{exercise}

\begin{exercise}[2]
Find the series for

\[
e^{-x^2}.
\]
\end{exercise}

\begin{exercise}[2]
Find the first four nonzero terms of

\[
\sqrt{1+x}.
\]
\end{exercise}

\begin{exercise}[2]
Find the first four nonzero terms of

\[
(1+x)^{-1/2}.
\]
\end{exercise}

\begin{exercise}[2]
Determine the radius of convergence of

\[
\sum_{n=0}^{\infty}
\frac{x^n}{n!}.
\]
\end{exercise}

\begin{exercise}[2]
Determine the radius of convergence of

\[
\sum_{n=1}^{\infty}
\frac{x^n}{n}.
\]
\end{exercise}

\begin{exercise}[3]
Determine the interval of convergence of

\[
\sum_{n=1}^{\infty}
\frac{x^n}{n}.
\]
\end{exercise}

\begin{exercise}[3]
Find the Taylor series of \(e^x\) about \(x=1\).
\end{exercise}

\begin{exercise}[3]
Find the Taylor series of \(\ln x\) about \(x=1\).
\end{exercise}

\begin{exercise}[3]
Use a series to evaluate

\[
\lim_{x\to0}
\frac{
e^x-1-x
}{
x^2
}.
\]
\end{exercise}

\begin{exercise}[3]
Use a series to evaluate

\[
\lim_{x\to0}
\frac{
\sin x-x
}{
x^3
}.
\]
\end{exercise}

\begin{exercise}[3]
Use a series to evaluate

\[
\lim_{x\to0}
\frac{
\tan x-x
}{
x^3
}.
\]
\end{exercise}

\begin{exercise}[3]
Use the binomial series to approximate

\[
\sqrt{1.01}.
\]
\end{exercise}

\begin{exercise}[3]
Use the arctangent series to approximate \(\pi\).
\end{exercise}

\begin{exercise}[3]
Derive Euler's formula from exponential, sine, and cosine series.
\end{exercise}

\begin{exercise}[3]
Compute the Cauchy product of two geometric series.
\end{exercise}

\begin{exercise}[3]
Find the ordinary generating function for

\[
a_n=2^n.
\]
\end{exercise}

\begin{exercise}[3]
Find the generating function for the sequence

\[
1,2,3,4,\ldots.
\]
\end{exercise}

\begin{exercise}[3]
Compute Fourier coefficients for an even function on \([-\pi,\pi]\).
\end{exercise}

\begin{exercise}[3]
Compute Fourier coefficients for an odd function on \([-\pi,\pi]\).
\end{exercise}

\begin{exercise}[3]
Write a Fourier series using complex exponentials.
\end{exercise}

\begin{exercise}[3]
Verify Parseval's identity for a simple trigonometric polynomial.
\end{exercise}

\begin{exercise}[4]
Prove that a power series can be differentiated termwise inside its radius of
convergence.
\end{exercise}

\begin{exercise}[4]
Prove that termwise integration preserves the radius of convergence.
\end{exercise}

\begin{exercise}[4]
Derive the Cauchy--Hadamard formula.
\end{exercise}

\begin{exercise}[4]
Derive the arcsine series from the binomial expansion.
\end{exercise}

\begin{exercise}[4]
Study a smooth function that is not equal to its Taylor series near the
expansion point.
\end{exercise}

\begin{exercise}[4]
Derive the first correction terms in Stirling's asymptotic expansion.
\end{exercise}

\begin{exercise}[4]
Study the Gibbs phenomenon for Fourier series near a jump discontinuity.
\end{exercise}

%%%%%%%%%%%%%%%%%%%%%%%%%%%%%%%%%%%%%%%%%%%%%%%%%%%%%%%%%%%%
\subsection{Conceptual Summary}
\label{subsec:series-expansions-summary}
%%%%%%%%%%%%%%%%%%%%%%%%%%%%%%%%%%%%%%%%%%%%%%%%%%%%%%%%%%%%

A power series has form

\[
\boxed{
\sum_{n=0}^{\infty}
c_n(x-a)^n.
}
\]

Taylor series encode derivatives at a point:

\[
\boxed{
f(x)
=
\sum_{n=0}^{\infty}
\frac{
f^{(n)}(a)
}{
n!
}
(x-a)^n
}
\]

when the series converges to \(f\).

Essential expansions include

\[
\boxed{
e^x
=
\sum_{n=0}^{\infty}
\frac{x^n}{n!},
}
\]

\[
\boxed{
\sin x
=
\sum_{n=0}^{\infty}
(-1)^n
\frac{x^{2n+1}}{(2n+1)!},
}
\]

\[
\boxed{
\cos x
=
\sum_{n=0}^{\infty}
(-1)^n
\frac{x^{2n}}{(2n)!},
}
\]

and

\[
\boxed{
\frac{1}{1-x}
=
\sum_{n=0}^{\infty}x^n,
\qquad
|x|<1.
}
\]

Series can be differentiated and integrated termwise inside their radius of
convergence.

Asymptotic expansions provide approximate descriptions even when the
corresponding infinite series does not converge.

Fourier series expand periodic functions in trigonometric or complex
exponential bases.

The central practical principles are

\[
\boxed{
\text{track the convergence region},
}
\]

\[
\boxed{
\text{estimate truncation error},
}
\]

and

\[
\boxed{
\text{distinguish exact series representations from asymptotic expansions}.
}
\]

%%%%%%%%%%%%%%%%%%%%%%%%%%%%%%%%%%%%%%%%%%%%%%%%%%%%%%%%%%%%
\subsection{Section Connections}
\label{subsec:series-expansions-connections}
%%%%%%%%%%%%%%%%%%%%%%%%%%%%%%%%%%%%%%%%%%%%%%%%%%%%%%%%%%%%

\chapterconnection{
Series Expansions connects local calculus with global approximation.
Taylor series build directly from the derivative formulas in
Section~\ref{sec:derivative-tables}; termwise integration connects with
Section~\ref{sec:integral-tables}; asymptotic expansions connect with Limit
Laws and Numerical Analysis; generating functions connect with Combinatorics;
Fourier series connect with Analysis, Differential Equations, Mathematical
Physics, and signal models. The next reference section, Transform Tables,
collects Laplace, Fourier, inverse-transform, convolution, shift, scaling,
differentiation, and selected \(z\)-transform formulas.
}

%%%%%%%%%%%%%%%%%%%%%%%%%%%%%%%%%%%%%%%%%%%%%%%%%%%%%%%%%%%%
% SECTION SOURCE TRACKING
%%%%%%%%%%%%%%%%%%%%%%%%%%%%%%%%%%%%%%%%%%%%%%%%%%%%%%%%%%%%

%------------------------------------------------------------
% 14.11 Series Expansions
%------------------------------------------------------------
%
% Source basis:
%   - Standard calculus and real/complex analysis.
%   - Standard power-series, Taylor-series, Fourier-series,
%     generating-function, and asymptotic formulas.
%
% Used for:
%   - sequences and series;
%   - geometric series;
%   - power series;
%   - radii and intervals of convergence;
%   - Taylor and Maclaurin series;
%   - Taylor remainders;
%   - exponential series;
%   - sine and cosine series;
%   - hyperbolic series;
%   - logarithm series;
%   - inverse-trigonometric series;
%   - binomial series;
%   - termwise differentiation and integration;
%   - Cauchy products;
%   - composition;
%   - approximation and error;
%   - asymptotic expansions;
%   - Stirling expansion;
%   - harmonic-number expansion;
%   - Fourier series;
%   - Fourier coefficients;
%   - Parseval identity;
%   - generating functions;
%   - special-function series;
%   - common errors;
%   - applications and exercises.
%
% Notes:
%   - Transform Tables is developed next in Section 14.12.
%   - Probability Distributions follows in Section 14.13.
%
%%%%%%%%%%%%%%%%%%%%%%%%%%%%%%%%%%%%%%%%%%%%%%%%%%%%%%%%%%%%